% from main.tex
%%%%%%%%%%  scan idx 259  |  folio 246  %%%%%%%%%%
\textbf{Proposition 3.6.} Soit
\[
\begin{tikzcd}
X \arrow[r,hook,"i"] \arrow[d,"f"'] & X' \arrow[dl,"f'"]\\
Y &
\end{tikzcd}
\]
un triangle commutatif de schémas, où \(i\) est une immersion fermée et \(f'\) un morphisme lisse, et soit \(E\in\Ob D(X)\). Pour que \(E\) soit localement d'amplitude plate finie rel. à \(f\), il faut et il suffit que \(i_*(E)\) soit localement d'amplitude plate finie (au sens absolu). Plus précisément, si \(f'\) est de dimension relative \(\leq d\), on a
\begin{align*}
\text{(i)}\quad &\tamp(i_*(E))\subset[a,b] \Longrightarrow \tamp_f(E)\subset[a,b],\\
\text{(ii)}\quad &\tamp_f(E)\subset[a,b] \Longrightarrow \tamp(i_*(E))\subset[a-d,b].
\end{align*}

